\documentclass[11pt,a4paper]{article}

% === 中文支持 ===
\usepackage{ctex}
\usepackage[T1]{fontenc}

% === 页面布局 ===
\usepackage[top=2.5cm, bottom=2.5cm, left=2.5cm, right=2.5cm]{geometry}

% === 表格 ===
\usepackage{booktabs}
\usepackage{longtable}
\usepackage{array}

% === 代码 ===
\usepackage{listings}
\lstset{
  basicstyle=\ttfamily\small,
  breaklines=true,
  frame=single,
  backgroundcolor=\color{gray!10},
  showstringspaces=false
}

% === 颜色 ===
\usepackage{xcolor}
\definecolor{primary}{HTML}{404040}
\definecolor{success}{HTML}{22c55e}
\definecolor{danger}{HTML}{e74c3c}
\definecolor{info}{HTML}{3b82f6}

% === 超链接 ===
\usepackage[hidelinks]{hyperref}
\hypersetup{
  colorlinks=true,
  linkcolor=primary,
  urlcolor=info,
  citecolor=primary
}

% === 标题样式 ===
\usepackage{titlesec}
\titleformat{\section}{\Large\bfseries\color{primary}}{}{0em}{}[\titlerule]
\titleformat{\subsection}{\large\bfseries\color{primary}}{}{0em}{}
\titleformat{\subsubsection}{\normalsize\bfseries\color{primary}}{}{0em}{}

% === 图形 ===
\usepackage{tikz}
\usetikzlibrary{shapes.geometric, arrows.meta, positioning, fit}

% === 文档信息 ===
\title{\textbf{AI Agent 应用商店}\\[0.5em]开源 Skill 架构调研报告}
\author{万能产品小助手（PM AI）\\魏子政 审阅}
\date{2026-05-06\\版本 v0.1}

\begin{document}

\maketitle

\begin{abstract}
本文档对企业级 AI Agent 应用商店项目的技术参考进行调研，重点分析开源 Skill 生态中与应用商店相关的架构设计。调研对象为 \texttt{npx skills}（skills.sh）平台上的应用商店类 Skill，包括 Flow Nexus App Store 和 Membrane Marketplacer。通过蒸馏其架构设计思路，为本项目提供参考。

\textbf{关键词}：AI Agent、应用商店、Skill、MCP、Marketplace、Flow Nexus、Membrane
\end{abstract}

\tableofcontents
\newpage

% ============================================================
\section{调研背景与目标}
% ============================================================

\subsection{调研背景}
企业级 AI Agent 应用商店项目（以下简称"本项目"）已进入 V0 准备阶段，核心目标是构建一个支持技能分发、安全扫描、审批和运营的应用商店平台。在正式开发前，需要对现有的开源 AI Agent 应用商店实现进行调研，了解其架构设计思路和最佳实践。

\subsection{调研目标}
\begin{enumerate}
  \item 了解现有 AI Agent 应用商店的架构模式
  \item 分析不同方案的优劣势
  \item 蒸馏可复用的设计思路，为本项目提供参考
  \item 评估哪些设计可以直接借鉴，哪些需要定制化
\end{enumerate}

\subsection{调研方法}
通过 \texttt{npx skills find} 命令搜索 skills.sh 平台上的应用商店相关 Skill，筛选出最相关的实现进行深入分析。共安装并分析了 3 个 Skill 包：
\begin{itemize}
  \item \texttt{ruvnet/ruflo@agent-app-store}（221 安装量）— Flow Nexus App Store
  \item \texttt{membranedev/application-skills@marketplacer}（29 安装量）— Membrane Marketplacer
  \item \texttt{aiskillstore/marketplace@webapp-testing}（174 安装量）— WebApp Testing（非应用商店，排除）
\end{itemize}

% ============================================================
\section{参考方案一：Flow Nexus App Store}
% ============================================================

\subsection{基本信息}
\begin{table}[h]
\centering
\begin{tabular}{ll}
\toprule
\textbf{属性} & \textbf{值} \\
\midrule
来源 & \texttt{ruvnet/ruflo@agent-app-store} \\
平台 & Flow Nexus \\
安装量 & 221 \\
安全评级 & Safe（Gen）/ Low Risk（Snyk） \\
适用范围 & Claude Code, OpenClaw, Cursor, Codex 等 \\
\bottomrule
\end{tabular}
\caption{Flow Nexus App Store 基本信息}
\end{table}

\subsection{架构概览}
Flow Nexus App Store 采用 \textbf{MCP（Model Context Protocol）工具调用架构}。Agent 通过调用预定义的 MCP 工具与平台交互，而非直接操作 API 或数据库。

核心架构组件：
\begin{enumerate}
  \item \textbf{MCP 工具层}：定义标准化的工具接口（app\_search、app\_store\_publish\_app、template\_deploy、app\_analytics）
  \item \textbf{Agent Skill 层}：封装业务逻辑和最佳实践，指导 Agent 如何使用工具
  \item \textbf{平台后端}：Flow Nexus 平台提供实际的存储、搜索、部署能力
\end{enumerate}

\subsection{核心功能分析}

\subsubsection{应用搜索与发现}
\begin{lstlisting}[caption={MCP 工具调用示例：搜索应用}]
mcp__flow-nexus__app_search({
  search: "authentication",
  category: "backend",
  featured: true,
  limit: 20
})
\end{lstlisting}

设计特点：
\begin{itemize}
  \item 支持关键词搜索 + 分类过滤 + 精选推荐
  \item AI 驱动的智能推荐（基于用户需求和历史）
  \item 8 大应用分类：Web APIs、Frontend、Full-Stack、CLI Tools、Data Processing、ML Models、Blockchain、Mobile
\end{itemize}

\subsubsection{应用发布}
\begin{lstlisting}[caption={MCP 工具调用示例：发布应用}]
mcp__flow-nexus__app_store_publish_app({
  name: "My Auth Service",
  description: "JWT-based authentication microservice",
  category: "backend",
  version: "1.0.0",
  source_code: sourceCode,
  tags: ["auth", "jwt", "express"]
})
\end{lstlisting}

发布流程包含：
\begin{itemize}
  \item 源代码直接提交（支持多种格式）
  \item 语义化版本管理（SemVer）
  \item 标签体系辅助发现
\end{itemize}

\subsubsection{模板部署}
\begin{lstlisting}[caption={MCP 工具调用示例：部署模板}]
mcp__flow-nexus__template_deploy({
  template_name: "express-api-starter",
  deployment_name: "my-api",
  variables: {
    api_key: "key",
    database_url: "postgres://..."
  }
})
\end{lstlisting}

设计特点：
\begin{itemize}
  \item 一键部署模板化应用
  \item 变量注入机制支持自定义配置
  \item 部署与发布分离（模板独立管理）
\end{itemize}

\subsubsection{质量保障体系}
Flow Nexus 定义了完整的质量标准：
\begin{itemize}
  \item \textbf{文档要求}：必须有完整的安装和使用说明
  \item \textbf{安全扫描}：所有发布应用必须通过漏洞评估
  \item \textbf{性能基准}：资源使用优化和性能测试
  \item \textbf{版本兼容}：向后兼容性管理
  \item \textbf{评分系统}：用户评分和评论机制
  \item \textbf{收益透明}：公平的收益分配策略
\end{itemize}

\subsection{架构图}
\begin{figure}[h]
\centering
\begin{tikzpicture}[
  node distance=1.5cm,
  box/.style={rectangle, draw, rounded corners, minimum width=3cm, minimum height=1cm, align=center, fill=blue!5},
  tool/.style={rectangle, draw, rounded corners, minimum width=2.5cm, minimum height=0.8cm, align=center, fill=green!10},
  arrow/.style={-{Stealth}, thick}
]
  \node[box] (agent) {AI Agent\\(Cursor/Claude)};
  \node[box, right=2cm of agent] (skill) {Skill 层\\(SKILL.md)};
  \node[box, right=2cm of skill] (mcp) {MCP 工具层\\(工具接口)};
  \node[box, below=1.5cm of mcp] (platform) {Flow Nexus 平台\\(后端服务)};

  \draw[arrow] (agent) -- (skill);
  \draw[arrow] (skill) -- (mcp);
  \draw[arrow] (mcp) -- (platform);

  \node[tool, above=0.5cm of mcp] (t1) {app\_search};
  \node[tool, left=0.3cm of t1] (t2) {app\_publish};
  \node[tool, right=0.3cm of t1] (t3) {template\_deploy};

  \draw[arrow] (t1) -- (mcp);
  \draw[arrow] (t2) -- (mcp);
  \draw[arrow] (t3) -- (mcp);
\end{tikzpicture}
\caption{Flow Nexus App Store 架构图}
\end{figure}

\subsection{优势与局限}
\begin{table}[h]
\centering
\begin{tabular}{p{5cm}p{8cm}}
\toprule
\textbf{优势} & \textbf{局限} \\
\midrule
MCP 工具标准化，易于扩展 & 依赖 Flow Nexus 平台后端 \\
Skill 层封装最佳实践 & 安全扫描细节未公开 \\
多 Agent 适配（Cursor/Claude/OpenClaw） & 缺少审批流程设计 \\
分类体系完整（8 大类） & 收益分配机制对私有化场景不适用 \\
模板化部署降低使用门槛 & 无企业级 RBAC/权限设计 \\
\bottomrule
\end{tabular}
\caption{Flow Nexus App Store 优劣势分析}
\end{table}

% ============================================================
\section{参考方案二：Membrane Marketplacer}
% ============================================================

\subsection{基本信息}
\begin{table}[h]
\centering
\begin{tabular}{ll}
\toprule
\textbf{属性} & \textbf{值} \\
\midrule
来源 & \texttt{membranedev/application-skills@marketplacer} \\
平台 & Membrane + Marketplacer \\
安装量 & 29 \\
安全评级 & Safe（Gen）/ Med Risk（Snyk） \\
协议 & MIT \\
\bottomrule
\end{tabular}
\caption{Membrane Marketplacer 基本信息}
\end{table}

\subsection{架构概览}
Membrane Marketplacer 采用 \textbf{中间件代理架构}。通过 Membrane CLI 作为统一入口，Agent 无需直接处理认证和 API 调用细节。

核心架构组件：
\begin{enumerate}
  \item \textbf{Membrane CLI}：统一的命令行工具，处理认证、连接管理、API 代理
  \item \textbf{Connection 层}：自动管理 OAuth 认证、凭证刷新、连接状态轮询
  \item \textbf{Action 层}：预构建的 API 操作（CRUD），支持自然语言搜索
  \item \textbf{Proxy 层}：当预构建 Action 不够时，直接代理 HTTP 请求
\end{enumerate}

\subsection{核心功能分析}

\subsubsection{认证与连接管理}
Membrane 的认证流程设计精巧：
\begin{itemize}
  \item \textbf{免密钥设计}：Agent 不需要持有 API Key，Membrane 服务端管理凭证
  \item \textbf{自动刷新}：Token 过期时透明刷新，Agent 无感知
  \item \textbf{连接状态机}：BUILDING -- READY / CLIENT\_ACTION\_REQUIRED / CONFIGURATION\_ERROR
  \item \textbf{多 Agent 适配}：支持 claude、openclaw、codex、warp、windsurf 等
\end{itemize}

\subsubsection{数据模型}
Marketplacer 定义了极其丰富的数据模型（40+ 实体）：

\begin{table}[h]
\centering
\small
\begin{tabular}{ll}
\toprule
\textbf{领域} & \textbf{实体} \\
\midrule
商品管理 & Listing、Listing Draft、Category、Brand、Attribute \\
交易流程 & Order、Transaction、Payment、Shipping Quote、Label \\
用户互动 & User、Review、Inquiry、Conversation、Message \\
售后保障 & Return Request、Return Reason、Dispute、Dispute Reason \\
运营管理 & Site、Theme、Template、Setting、Announcement \\
财务系统 & Plan、Invoice、Coupon、Credit \\
系统基础 & Authentication、Authorization、Webhook、Role、Permission \\
国际化 & Currency、Language、Country、Region、Tax Rate \\
\bottomrule
\end{tabular}
\caption{Marketplacer 数据模型概览}
\end{table}

\subsubsection{Action 发现与执行}
\begin{lstlisting}[caption={Action 发现示例}]
# 自然语言搜索可用操作
membrane action list --connectionId=CONN_ID \
  --intent "QUERY" --limit 10 --json

# 执行预构建操作
membrane action run ACTION_ID \
  --connectionId=CONN_ID --input '{"key": "value"}' --json

# 直接代理 API 请求（当预构建操作不够时）
membrane request CONN_ID /path/to/endpoint \
  -X POST --json -d '{"key": "value"}'
\end{lstlisting}

\subsection{架构图}
\begin{figure}[h]
\centering
\begin{tikzpicture}[
  node distance=1.5cm,
  box/.style={rectangle, draw, rounded corners, minimum width=3cm, minimum height=1cm, align=center, fill=orange!5},
  arrow/.style={-{Stealth}, thick}
]
  \node[box] (agent) {AI Agent\\(任意 CLI Agent)};
  \node[box, right=2cm of agent] (cli) {Membrane CLI\\(统一入口)};
  \node[box, right=2cm of cli] (conn) {Connection 层\\(认证 + 状态机)};
  \node[box, below left=1.5cm and 0.5cm of conn] (action) {Action 层\\(预构建操作)};
  \node[box, below right=1.5cm and 0.5cm of conn] (proxy) {Proxy 层\\(直接 API 代理)};
  \node[box, below=1.5cm of conn] (api) {Marketplacer API\\(第三方平台)};

  \draw[arrow] (agent) -- (cli);
  \draw[arrow] (cli) -- (conn);
  \draw[arrow] (conn) -- (action);
  \draw[arrow] (conn) -- (proxy);
  \draw[arrow] (action) -- (api);
  \draw[arrow] (proxy) -- (api);
\end{tikzpicture}
\caption{Membrane Marketplacer 架构图}
\end{figure}

\subsection{优势与局限}
\begin{table}[h]
\centering
\begin{tabular}{p{5cm}p{8cm}}
\toprule
\textbf{优势} & \textbf{局限} \\
\midrule
免密钥设计，安全性高 & 依赖外部 Membrane 服务 \\
连接状态机设计完善 & Med Risk 安全评级 \\
自然语言 Action 发现 & 主要面向电商场景（非 AI Skill） \\
数据模型极其丰富（40+ 实体） & 本身是集成工具而非独立应用商店 \\
MIT 开源协议 & 不适合直接作为企业私有化方案 \\
\bottomrule
\end{tabular}
\caption{Membrane Marketplacer 优劣势分析}
\end{table}

% ============================================================
\section{对比分析}
% ============================================================

\begin{table}[h]
\centering
\small
\begin{tabular}{p{2.5cm}p{4.5cm}p{4.5cm}}
\toprule
\textbf{维度} & \textbf{Flow Nexus} & \textbf{Membrane} \\
\midrule
架构模式 & MCP 工具调用 & 中间件代理 \\
核心定位 & AI Skill 应用商店 & 通用 Marketplace 集成 \\
认证方式 & 平台内建 & 免密钥代理 \\
安全扫描 & 有（漏洞评估） & 无（依赖第三方） \\
审批流程 & 无 & 无 \\
部署能力 & 模板化一键部署 & 无 \\
数据模型 & 简洁（应用+模板+分析） & 极丰富（40+ 实体） \\
开源协议 & 未声明 & MIT \\
私有化适配 & 低（依赖平台） & 低（依赖 Membrane 服务） \\
企业级 RBAC & 无 & 有（Role/Permission） \\
\bottomrule
\end{tabular}
\caption{两个方案的全面对比}
\end{table}

% ============================================================
\section{对本项目的启示}
% ============================================================

\subsection{可直接借鉴的设计}

\begin{enumerate}
  \item \textbf{分类体系}：Flow Nexus 的 8 大分类（Web APIs、Frontend、Full-Stack、CLI Tools、Data Processing、ML Models、Blockchain、Mobile）可作为本项目分类设计的参考基线。
  
  \item \textbf{质量保障标准}：文档要求 + 安全扫描 + 性能基准 + 版本兼容 + 评分系统的组合思路值得借鉴，但需要根据企业场景定制。
  
  \item \textbf{模板化部署}：变量注入 + 一键部署的设计思路适用于企业内部技能的分发场景。
  
  \item \textbf{免密钥代理思路}：Membrane 的 Connection 状态机设计可用于本项目与外部 API 集成时的认证管理。
  
  \item \textbf{数据模型丰富度}：Marketplacer 的 40+ 实体模型展示了企业级应用商店可能需要的数据维度，本项目可按需裁剪。
\end{enumerate}

\subsection{本项目需要额外设计的部分}

\begin{enumerate}
  \item \textbf{审批工作流}：两个方案均无审批流程，但企业场景必须。本项目已规划状态机方案。
  
  \item \textbf{三层安全扫描}：Semgrep + ClamAV + LLM 语义分析，两个方案的安全扫描均较浅。
  
  \item \textbf{企业级 RBAC}：本项目已规划 4 角色权限体系（admin/reviewer/developer/viewer）。
  
  \item \textbf{私有化部署}：两个方案均依赖外部平台，本项目需完全私有化。
  
  \item \textbf{审计追溯}：企业场景的审计日志、操作追溯是两个方案均未覆盖的。
\end{enumerate}

\subsection{能力扩充方向}
基于调研结果，建议后续关注以下能力扩充方向：
\begin{itemize}
  \item \textbf{MCP 协议集成}：考虑将本项目的 API 封装为 MCP 工具，使外部 AI Agent 可直接调用
  \item \textbf{Skill CLI 生态对接}：支持通过 \texttt{npx skills} 发布和发现本项目的技能包
  \item \textbf{Figma MCP 集成}：如果魏子政提供设计稿，可评估 Figma MCP 用于设计令牌同步
\end{itemize}

% ============================================================
\section{结论}
% ============================================================

现有开源 AI Agent 应用商店方案处于早期阶段，两个主要参考方案各有侧重但均不完整：
\begin{itemize}
  \item Flow Nexus 提供了较好的 Skill 管理和分发框架，但缺乏企业级安全治理和审批流程
  \item Membrane 提供了优秀的认证管理和数据模型设计，但本身是集成工具而非独立应用商店
\end{itemize}

本项目的差异化价值在于：
\begin{enumerate}
  \item \textbf{企业级安全治理}：三层扫描 + 审批工作流 + 审计追溯
  \item \textbf{完全私有化}：不依赖任何外部平台
  \item \textbf{双智能体协作}：OpenClaw（PM）+ Cursor（开发）的明确分工
  \item \textbf{规范化文档体系}：V0 阶段即建立完整的文档生命周期管理
\end{enumerate}

调研结论：现有开源方案可作为参考，但本项目的架构需要自研，不依赖任何单一开源方案。

\end{document}
